\documentclass[10pt,a4paper]{amsart}
\usepackage{amsmath,amssymb,amsthm,graphicx}
% Retained preamble for the Stacks Project chapter snapshot in this directory.
% The chapter's own preamble.tex is a copy of the shared Stacks Project preamble; the block
% below is sliced MECHANICALLY, line for line, from it. Nothing was added, removed or
% redefined; the wrapper class, the xr-hyper/hyperref wiring and the externaldocument list
% are omitted because the assembler supplies the document class.
%
% lines taken from preamble.tex: 3, 5, 9, 11, 12, 156, 157, 158, ... (51 lines total)
%
% package coverage restored from the chapter's own preamble (2026-10-13): xy, xy-2cell

\usepackage[all]{xy}
\xyoption{2cell}
\UseAllTwocells

\newenvironment{reference}{\comment}{\endcomment}
\newenvironment{slogan}{\comment}{\endcomment}
\newenvironment{history}{\comment}{\endcomment}
\theoremstyle{plain}
\newtheorem{theorem}[subsection]{Theorem}
\newtheorem{proposition}[subsection]{Proposition}
\newtheorem{lemma}[subsection]{Lemma}
\theoremstyle{definition}
\newtheorem{definition}[subsection]{Definition}
\newtheorem{example}[subsection]{Example}
\newtheorem{exercise}[subsection]{Exercise}
\newtheorem{situation}[subsection]{Situation}
\theoremstyle{remark}
\newtheorem{remark}[subsection]{Remark}
\newtheorem{remarks}[subsection]{Remarks}
\numberwithin{equation}{subsection}
\def\lim{\mathop{\mathrm{lim}}\nolimits}
\def\colim{\mathop{\mathrm{colim}}\nolimits}
\def\Spec{\mathop{\mathrm{Spec}}}
\def\Hom{\mathop{\mathrm{Hom}}\nolimits}
\def\Ext{\mathop{\mathrm{Ext}}\nolimits}
\def\SheafHom{\mathop{\mathcal{H}\!\mathit{om}}\nolimits}
\def\SheafExt{\mathop{\mathcal{E}\!\mathit{xt}}\nolimits}
\def\Sch{\mathit{Sch}}
\def\Mor{\mathop{\mathrm{Mor}}\nolimits}
\def\Ob{\mathop{\mathrm{Ob}}\nolimits}
\def\Sh{\mathop{\mathit{Sh}}\nolimits}
\def\NL{\mathop{N\!L}\nolimits}
\def\CH{\mathop{\mathrm{CH}}\nolimits}
\def\proetale{{pro\text{-}\acute{e}tale}}
\def\etale{{\acute{e}tale}}
\def\QCoh{\mathit{QCoh}}
\def\Ker{\mathop{\mathrm{Ker}}}
\def\Im{\mathop{\mathrm{Im}}}
\def\Coker{\mathop{\mathrm{Coker}}}
\def\Coim{\mathop{\mathrm{Coim}}}
\DeclareMathSymbol{\boxtimes}{\mathbin}{AMSa}{"02}
\def\QCohstack{\mathcal{QC}\!\mathit{oh}}
\def\Cohstack{\mathcal{C}\!\mathit{oh}}
\def\Spacesstack{\mathcal{S}\!\mathit{paces}}
\def\Quotfunctor{\mathrm{Quot}}
\def\Hilbfunctor{\mathrm{Hilb}}
\def\Curvesstack{\mathcal{C}\!\mathit{urves}}
\def\Polarizedstack{\mathcal{P}\!\mathit{olarized}}
\def\Complexesstack{\mathcal{C}\!\mathit{omplexes}}
\def\Pic{\mathop{\mathrm{Pic}}\nolimits}
\def\Picardstack{\mathcal{P}\!\mathit{ic}}
\def\Picardfunctor{\mathrm{Pic}}
\def\Deformationcategory{\mathcal{D}\!\mathit{ef}}
\usepackage{cleveref}
\crefname{originallemma}{original lemma}{original lemmas}
\Crefname{originallemma}{Original lemma}{Original lemmas}
\makeatletter
\newlabel{algebraic-lemma-stack-presentation}{{04T5}{source}{}{sourcepointer-1}{}}
\newlabel{algebraic-lemma-stack-presentation@cref}{{[originallemma][0][]04T5}{[1][0][]source}{}{}{}}
\makeatother
\crefname{originallemma}{original lemma}{original lemmas}
\Crefname{originallemma}{Original lemma}{Original lemmas}
\makeatletter
\newlabel{bootstrap-lemma-quotient-stack-isom}{{04TB}{source}{}{sourcepointer-2}{}}
\newlabel{bootstrap-lemma-quotient-stack-isom@cref}{{[originallemma][0][]04TB}{[1][0][]source}{}{}{}}
\makeatother
\crefname{originallemma}{original lemma}{original lemmas}
\Crefname{originallemma}{Original lemma}{Original lemmas}
\makeatletter
\newlabel{categories-lemma-fibred-groupoids-fibre-product-goes-up}{{06N6}{source}{}{sourcepointer-3}{}}
\newlabel{categories-lemma-fibred-groupoids-fibre-product-goes-up@cref}{{[originallemma][0][]06N6}{[1][0][]source}{}{}{}}
\makeatother
\crefname{originallemma}{original lemma}{original lemmas}
\Crefname{originallemma}{Original lemma}{Original lemmas}
\makeatletter
\newlabel{spaces-morphisms-lemma-section-immersion}{{03KP}{source}{}{sourcepointer-4}{}}
\newlabel{spaces-morphisms-lemma-section-immersion@cref}{{[originallemma][0][]03KP}{[1][0][]source}{}{}{}}
\makeatother
\crefname{originallemma}{original lemma}{original lemmas}
\Crefname{originallemma}{Original lemma}{Original lemmas}
\makeatletter
\newlabel{stacks-definition-mor-presheaf}{{02ZB}{source}{}{sourcepointer-5}{}}
\newlabel{stacks-definition-mor-presheaf@cref}{{[originallemma][0][]02ZB}{[1][0][]source}{}{}{}}
\makeatother
\title{Stacks, Sheaves on Algebraic Stacks: lemma-fibre-products}
\author{The Stacks Project Authors}
\date{stacks-project:stacks-sheaves.tex:local-snapshot-20261012; GNU FDL 1.2 or later}
\begin{document}
\maketitle
\noindent\textit{Editorial note.} Statement span is lines 2863-2874 of sources/stacks-algebraic-stacks-sheaves/pinned-original-source.tex; proof span is lines 2876-2911, the complete author proof (34 body lines).\par\medskip
\noindent\textit{Editorial note.} Navigation stubs: 5 cited-but-not-included labels are supplied as non-load-bearing pointers with their pinned tags where the snapshot tag index has them.\par\medskip
\noindent\textit{Editorial note.} extra\_packages is empty on purpose: the retained preamble already loads xy with options and a second declaration clashes in the assembler wrapper.\par\medskip
\noindent\textit{Editorial note.} Text is retained verbatim, including the author's own typos and placeholders.\par\medskip
\noindent\textit{Editorial note.} Queue index 13 of queue\_algebraic-stacks-sheaves.json (pb=1841, stmt 2863-2874, proof 2876-2911).\par\medskip
\noindent\textit{Editorial note.} Cross-chapter same-label warning (NOT an occupancy conflict, Chief ruling 2026-10-13 evening): label 'lemma-fibre-products' also exists in stacks-schemes — a different proposition in a different chapter file.\par\medskip
\noindent\textit{Allowed standard prerequisites (editorial).}\begin{itemize}
\item Explicit prerequisite boundary. The cited inputs are prerequisites; the argument in the retained span is not replaced by any cited result.
\end{itemize}\medskip
\ifdefined\hypertarget\hypertarget{sourcepointer-1}{}\fi
\noindent\textit{Source reference.} Original lemma 04T5: sources/stacks-algebraic-stacks-sheaves/pinned-original-source.tex, source lines 2911--2911. Cited by the retained statement or proof but not part of the collected spans; supplied only as a navigation pointer. The argument does not depend on its content.\par\medskip
\ifdefined\hypertarget\hypertarget{sourcepointer-2}{}\fi
\noindent\textit{Source reference.} Original lemma 04TB: sources/stacks-algebraic-stacks-sheaves/pinned-original-source.tex, source lines 2911--2911. Cited by the retained statement or proof but not part of the collected spans; supplied only as a navigation pointer. The argument does not depend on its content.\par\medskip
\ifdefined\hypertarget\hypertarget{sourcepointer-3}{}\fi
\noindent\textit{Source reference.} Original lemma 06N6: sources/stacks-algebraic-stacks-sheaves/pinned-original-source.tex, source lines 2911--2911. Cited by the retained statement or proof but not part of the collected spans; supplied only as a navigation pointer. The argument does not depend on its content.\par\medskip
\ifdefined\hypertarget\hypertarget{sourcepointer-4}{}\fi
\noindent\textit{Source reference.} Original lemma 03KP: sources/stacks-algebraic-stacks-sheaves/pinned-original-source.tex, source lines 2911--2911. Cited by the retained statement or proof but not part of the collected spans; supplied only as a navigation pointer. The argument does not depend on its content.\par\medskip
\ifdefined\hypertarget\hypertarget{sourcepointer-5}{}\fi
\noindent\textit{Source reference.} Original lemma 02ZB: sources/stacks-algebraic-stacks-sheaves/pinned-original-source.tex, source lines 2911--2911. Cited by the retained statement or proof but not part of the collected spans; supplied only as a navigation pointer. The argument does not depend on its content.\par\medskip
\begin{lemma}
\label{lemma-fibre-products}
Let $p : \mathcal{X} \to (\Sch/S)_{fppf}$
be a category fibred in groupoids.
\begin{enumerate}
\item The category $\mathcal{X}$ has fibre products.
\item If the $\mathit{Isom}$-presheaves of $\mathcal{X}$
are representable by algebraic spaces, then $\mathcal{X}$ has equalizers.
\item If $\mathcal{X}$ is an algebraic stack (or more generally
a quotient stack), then $\mathcal{X}$ has equalizers.
\end{enumerate}
\end{lemma}

\begin{proof}
Part (1) follows
Categories, Lemma \ref{categories-lemma-fibred-groupoids-fibre-product-goes-up}
as $(\Sch/S)_{fppf}$ has fibre products.

\medskip\noindent
Let $a, b : x \to y$ be morphisms of $\mathcal{X}$.
Set $U = p(x)$ and $V = p(y)$. The category of schemes has equalizers
hence we can let $W \to U$ be the equalizer of $p(a)$ and $p(b)$.
Denote $c : z \to x$ a morphism of $\mathcal{X}$ lying over $W \to U$.
The equalizer of $a$ and $b$, if it exists, is the equalizer of $a \circ c$
and $b \circ c$. Thus we may assume that $p(a) = p(b) = f : U \to V$.
As $\mathcal{X}$ is fibred in groupoids, there exists a unique automorphism
$i : x \to x$ in the fibre category of $\mathcal{X}$ over $U$ such that
$a \circ i = b$. Again the equalizer of $a$ and $b$ is the equalizer
of $\text{id}_x$ and $i$. Recall that the $\mathit{Isom}_\mathcal{X}(x)$
is the presheaf on $(\Sch/U)_{fppf}$ which to
$T/U$ associates the set of automorphisms of $x|_T$ in the fibre category
of $\mathcal{X}$ over $T$, see
Stacks, Definition \ref{stacks-definition-mor-presheaf}.
If $\mathit{Isom}_\mathcal{X}(x)$ is representable by an algebraic space
$G \to U$, then we see that $\text{id}_x$ and $i$ define morphisms
$e, i : U \to G$ over $U$. Set $M = U \times_{e, G, i} U$, which by
Morphisms of Spaces, Lemma \ref{spaces-morphisms-lemma-section-immersion}
is a scheme. Then it is clear that $x|_M \to x$ is the equalizer of
the maps $\text{id}_x$ and $i$ in $\mathcal{X}$.
This proves (2).

\medskip\noindent
If $\mathcal{X} = [U/R]$ for some groupoid in algebraic spaces
$(U, R, s, t, c)$ over $S$, then the hypothesis of (2) holds by
Bootstrap, Lemma \ref{bootstrap-lemma-quotient-stack-isom}.
If $\mathcal{X}$ is an algebraic stack, then we can choose a
presentation $[U/R] \cong \mathcal{X}$ by
Algebraic Stacks, Lemma \ref{algebraic-lemma-stack-presentation}.
\end{proof}

\end{document}
